\section{Related Work}

\subsection{LLM-as-Judge and Multi-Judge Evaluation}

LLM judges gained traction after \citet{zheng2023judging} established them as scalable alternatives to human annotation. \citet{wang2023fair} documented systematic failure modes (position bias, verbosity, sycophancy) and \citet{ye2024justice} provided a comprehensive bias taxonomy, motivating multi-judge panels, while \citet{raina2024judgerobust} and \citet{maloyan2025investigating} quantified adversarial vulnerabilities including prompt injection. Multi-judge aggregation extends the classical noisy annotator framework \citep{dawid1979maximum}: \citet{verga2024replacing} showed diverse panels can match single large judges, and \citet{panickssery2024llm} identified self-preference bias underlying within-family correlations. Most relevant, \citet{almasoud2026security} reported that 7-model committees reduce ASR from 50--68\% to 10--19\% but did not separate diversity from individual robustness, the gap we address. Concurrent work benchmarks judge quality \citep{tan2024judgebench,li2024autoj} and judge--human agreement across tasks \citep{bavaresco2025llms}, adversarial risks in multi-agent systems \citep{tamas2025benchmarking}, and collaborative evaluation under benign conditions \citep{qian2026collabeval}; none tests whether aggregation benefits survive adversarial manipulation.

\subsection{Fault Tolerance, Voting Theory, and Ensemble Diversity}

The Condorcet Jury Theorem (CJT) \citep{condorcet1785essai} guarantees majority-vote convergence when voters are independent and better than chance; \citet{ladha1992condorcet} showed positive correlation weakens this guarantee. Classical BFT \citep{lamport1982byzantine} tolerates non-strategic faults but not adversaries targeting system structure; our framework adopts the Stackelberg model of \citet{tambe2011security} with the weakest-link principle \citep{schneier2000secrets}. Correlated LLM judge errors \citep{kim2025correlated,kuai2026independent} violate both CJT and BFT independence assumptions; \citet{kuai2026independent} propose DBEI (difficulty-weighted entanglement), while $K_{\mathrm{eff}}$ captures a complementary dimension via pairwise MI. Mitigation strategies include Bayesian calibration and credibility scoring \citep{ebrahimi2025adversary}.

In ensemble theory, \citet{krogh1995neural} showed error decreases with diversity, but \citet{brown2005managing} demonstrated this holds only when members maintain low bias, a condition violated when LLM panels include weak judges for family coverage. \citet{kuncheva2003measures} found no pairwise diversity measure reliably predicts improvement, which we replicate: $K_{\mathrm{eff}}$ and $\kappa$-diversity yield opposite $K$-dependent trajectories. Our positive $K_{\mathrm{eff}}$--ASR association under adversarial conditions instantiates Brown et al.'s caveat: diversity co-occurs with capability reduction, and the cost dominates. The position-bias exception (diversity \emph{reduces} vulnerability for shared-bias attacks) and the operationalization-dependent trajectory do not follow from classical theory. Recent frameworks \citep{yang2026diversity,turkmen2026ensemble} characterize diversity under benign evaluation; we provide the adversarial complement.

\subsection{Adversarial Robustness of Evaluation Systems}

Adversarial vulnerability \citep{goodfellow2015adversarial} and transferability \citep{papernot2016transferability} are well-studied in vision, where controlled diversity via uncorrelated gradients \citep{kariyappa2019improving}, diversity regularization \citep{pang2019improving}, or vulnerability diversification \citep{zhu2020dverge} strengthens ensemble robustness, but requires coordinating member training, a control absent in post-hoc LLM judge selection. In NLP, \citet{wallace2019universal} and \citet{zou2023universal} showed adversarial triggers transfer across models. Our setting differs: adversaries use natural-language manipulation rather than gradient perturbations, aggregation rules are a controllable design parameter, and the 15-judge combinatorial design induces correlated residuals addressed via wild bootstrap \citep{mackinnon2018wild} with max-$p$ correction (\S\ref{sec:robustness_method}). Figure~\ref{fig:theory_comparison} (Appendix~\ref{app:theory_comparison}) summarizes the three assumption violations.
